\begin{tikzpicture}[font=\small,line cap=round,line join=round]\draw[geoTrait](0,0)--(3.6,0)--(2.4,1.8)--(0,1.8)--cycle;\node[below,font=\tiny]at(1.8,-.08){6 cm};\node[above,font=\tiny]at(1.2,1.82){4 cm};\node[left,font=\tiny]at(0,.9){3 cm};\draw[otNavy](.22,0)--(.22,.22)--(0,.22);\draw[otNavy](0,1.58)--(.22,1.58)--(.22,1.8);\geoSolution{\draw[geoAide](2.4,0)--(2.4,1.8);\draw[otGreenDark](2.4,.22)--(2.18,.22)(2.18,0)--(2.18,.22);\node[font=\tiny,otGreenDark]at(1.2,.9){rectangle};\node[font=\tiny,otGreenDark]at(2.95,.55){triangle};}\end{tikzpicture}\par\vspace{1.5mm}{\otRegular\fontsize{8}{10}\selectfont\color{otMuted}La perpendiculaire abaissée découpe le trapèze rectangle en un rectangle et un triangle rectangle.\par}
